% Transcription — fonds Alexandre Grothendieck, shelfmark 154, batch 6 (pages 101-120).
% Established from the facsimile archives/batches/154/batch-06.pdf, cut from a
% copy of 154.pdf fetched through the project's own relay
% (https://grothendieck-relay.onrender.com/source/154.pdf); PDF page k =
% archivists' page 100+k, no cover sheet. Per the skill transcribe-grothendieck.
% The folder is 175 pages in nine batches, transcribed in parallel; this is
% the sixth. Batch 5 (transcripts/154/batch-05.fr.tex) was read first for its
% notation; notation otherwise follows folder 80 (Systèmes de pseudo-droites).
%
% Pass: Opus 5.5 (claude-opus-5-5), 2026-09-26 — first pass, unchecked against the pages by a human.
%
% Dating and title. \dating is the folder's own date, copied verbatim from
% src/content/catalogue.ts; the folder belongs to the inventory group
% « Espaces stratifiés » (150-151), but since it has a date of its own the
% group's range is not added. Title from the catalogue, trailing full stop
% dropped; « [Système de pseudo-droites] » is bracketed there because the
% archivists supplied it.
%
% Paper. Nearly every leaf is fan-fold computer listing paper; the printed
% job banners (JESNEWS of 9 and 18 August 1982 on 101, 102, 111; job logs of
% 16 and 17 August 1982 on 109, 116, 120) are machine dates of the paper, not
% his. Pages 109, 111 and 120 have the print upside down relative to his
% drawings; read in the orientation of his hand. Page 110 is a portrait sheet.
%
% Pages skipped: none. Every leaf carries drawings of his.
% Pages summarised: none. Pages given only as a figure described in a French
% \note, with their few labels and formulas set: 101, 102, 105, 107, 108, 109,
% 114, 115, 116, 120; largely figures with some text or formulas: 110, 111,
% 112, 113, 118, 119.
% Pages transcribed from his hand: 103-104 (continuous), 106 (legends), 117
% (cancelled by a large X, read anyway), 119 (Euler characteristic count).
%
% His pagination and dates. No page carries a number or a date of his.
% Numbered runs: on p. 103 the list a) b) c) and the cases I 1) 2), II 1) 2);
% on p. 104 the cycle of positions labelled (1) (2) (3) (4) (1), with circled
% sector numbers 1, 2, 1', 2', 3', 4', « 3' = 1 », « 5' = 1 »; on p. 117 the
% alineas d)? e) f) (the letter of the first is uncertain; an a)-c) run must
% precede, presumably on a leaf outside this batch); on p. 119 an alinea b).
% A « Tantôt » / « cf. » cross-reference to an earlier page: « la
% rectification b) » on p. 117.
%
% What the batch is about.
%   103-104 The rotation operation rho(D, omega, a) = (D', omega', a') around
%           a face, driven by an oriented transition arc a on the double cover
%           D~ of a relative position D: a) D' from D by crossing a, b) omega'
%           opposite to the transport of omega, c) a' the transition arc on D~'
%           following the transport of a. Balayage / antibalayage of
%           specialisation arcs; eight cases in two pairs (a specialisation or
%           generisation arc); two ways to compute a'; a cycle of five
%           positions drawn; specialisation + generisation arcs = transition
%           arcs, vs. arcs d'isotopie; a count 2nu - 2 + (2+2) = 2nu + 2.
%   106     Legends on regions « générisée I » / « spécialisée II » and the
%           arc joining them; the blue dotted line of the cutting.
%   110     Counts on the supersingular band: n - nu, (n-1) - (nu-1) - (nu'-1),
%           2n - (nu + nu' - 1); a sign epsilon = +1 if the special face is
%           of vertex type (ultra-...), -1 otherwise; « how to compute eps? ».
%   111-112 Gluing near a vertex s: x vs x-bar not on the same line; the
%           vertex s would disappear (« non ! »); D_1 ∩ Delta_D' etc.; « je
%           n'arrive pas à suivre les "circuits autour d'un trou" ».
%   113     N = number of « green » vertices of F (vertices of order 4),
%           N' for F'; N' = N - 1 or N' = N; alpha', alpha'' >= 1, gamma >= 2.
%   117     Cancelled draft: a canonical corridor around K_0, isomorphic to a
%           tubular neighbourhood of D_i~ in X~; e) exceptional faces F*
%           adjacent to D_i give back the system Sigma; f) after rectification
%           b) exceptional lines of K are no longer characterised by containing
%           no vertex of L'; K'_0 as union of components of the unfolding K'
%           with no global boundary.
%   118     (X, K, (delta_s), (nu_s), ...); omega(D) -> omega(D') by
%           rotation / combinatorial transport; rho_D ≃ omega_D'.
%   119     chi(X) = s0 - s1 + s2 = 1 - g; vertices of order 4 so s1 = 2 s0,
%           g = 1 + s0 - s2, s2 = number of relative positions (supersingular
%           excluded), s0 = number of pairs (D, a).
%   101, 102, 105, 107-109, 114-116, 120  Drawings: arrangements in discs,
%           swept bands, figure-eight double covers with marked arcs, etc.
%
% Notation fixed once for the batch (as in batch 5): \mathcal{X} for his
% cursive X (surface of relative positions); \widetilde{D} for the double
% cover of a position D; \operatorname{Pol}(D); « p.r. » for position
% relative; \rho for his rotation operation.
%
% Readings to check first: the first line of p. 103 (« rotation » and
% « face »); p. 112 middle column; the opening two lines and the alinea d) of
% p. 117; the struck lines at the foot of p. 119; the first sign of the count
% on p. 104 (2 or >=).

\documentclass[11pt,a4paper]{article}
% Le préambule commun à toutes les transcriptions du fonds Grothendieck.
%
% Il tient en une idée : **l'incertitude se compose, elle ne se résout pas.**
% Une transcription qui lisse un mot douteux a détruit la seule chose pour
% laquelle on la faisait — un lecteur doit pouvoir savoir, à chaque endroit,
% ce qui était sur le feuillet et ce qui a été ajouté. Les six macros qui
% suivent sont donc l'appareil critique, et non de la décoration.
%
% Les mêmes commandes sont reconnues par `scripts/render.mjs`, qui produit la
% vue de lecture du site. Une macro ajoutée ici doit l'être aussi là-bas,
% faute de quoi le rendu échouera bruyamment — ce qui est voulu : un rendu
% silencieusement faux est pire qu'un rendu absent.

\ProvidesPackage{grothendieck}[2026/08/08 Transcriptions du fonds Grothendieck]

\usepackage{amsmath,amssymb,amsthm}
\usepackage{mathrsfs}
\usepackage{stmaryrd}
% Les diagrammes commutatifs sont partout dans ce fonds : c'est la langue
% même de la géométrie algébrique de Grothendieck, et une transcription qui
% les rendrait en prose ne transcrirait rien.
\usepackage{tikz-cd}
% Français par défaut — c'est la langue des feuillets — mais l'anglais est
% chargé aussi : la traduction et le résumé sont des documents anglais, et la
% typographie française (espaces avant les deux-points) y serait une faute.
% Ces documents basculent par \selectlanguage{english} en tête de corps.
\usepackage[english,french]{babel}
\usepackage{fontspec}
\usepackage{geometry}
\geometry{a4paper,margin=25mm,marginparwidth=28mm}
\usepackage{marginnote}
\usepackage{xcolor}
% ulem, not soul. `soul`'s \st and \ul are built on a character-by-character
% scan that XeLaTeX's font machinery breaks: the document fails with an
% undefined control sequence rather than degrading. ulem does the same two jobs
% and is Unicode-engine safe. `normalem` stops it hijacking \emph, which the
% transcriptions use for Grothendieck's own emphasis.
\usepackage[normalem]{ulem}
\usepackage{hyperref}
\hypersetup{colorlinks=true,linkcolor=black,urlcolor=black,pdfborder={0 0 0}}

% --- Métadonnées du lot -----------------------------------------------------
% Reprises telles quelles par le script de rendu, qui les affiche en tête de
% la vue de lecture. Elles fixent ce que le fichier prétend couvrir : sans
% elles, une transcription flotte sans point d'attache dans un fonds de seize
% mille feuillets.

\newcommand{\folder}[1]{\gdef\grothFolder{#1}}
\newcommand{\batch}[1]{\gdef\grothBatch{#1}}
\newcommand{\leaves}[2]{\gdef\grothFirst{#1}\gdef\grothLast{#2}}
\newcommand{\foldertitle}[1]{\gdef\grothFoldertitle{#1}}
\gdef\grothFolder{?}\gdef\grothBatch{?}\gdef\grothFirst{?}\gdef\grothLast{?}\gdef\grothFoldertitle{}

% --- Le résumé --------------------------------------------------------------
%
% Il ouvre la lecture modernisée, sous le titre « Résumé », et il est composé
% en petit. Ce n'est pas un ornement : c'est le seul passage du document qui
% ne soit pas des mathématiques, et un lecteur doit pouvoir le reconnaître
% avant de l'avoir lu — puis le sauter s'il connaît déjà le sujet, ou s'y
% arrêter s'il ne le connaît pas. Le filet qui le suit marque la reprise du
% corps.

\newenvironment{resume}
  {\small\setlength{\parskip}{0.45em}}
  {\par\medskip\hrule\medskip}

% --- Mots-clés ---------------------------------------------------------------
%
% En anglais, en clôture du résumé de la lecture modernisée : le vocabulaire
% moderne sous lequel chercher ce que les feuillets construisent. Le manifeste
% du site (`npm run manifest`) les extrait pour étiqueter la cote entière —
% cette ligne est la source unique des tags, il n'y a pas de fichier à côté.
\newcommand{\keywords}[1]{\par\smallskip\noindent
  {\footnotesize\textsc{Keywords} --- \textit{#1}}\par}

% --- L'appareil critique ----------------------------------------------------

% Le numéro de feuillet, en marge. C'est l'ancre qui permet de régler un
% désaccord sur une formule en regardant *un* feuillet plutôt qu'une chemise
% de six cents. Le site s'en sert aussi pour tourner les pages du fac-similé
% quand on fait défiler la transcription.
\newcommand{\leaf}[1]{%
  \marginnote{\footnotesize\color{gray!70}\textsf{#1}}%
  \label{leaf:#1}\ignorespaces}

% Illisible : on ne devine pas. Un mot inventé qui se lit comme les autres est
% le pire résultat possible.
\newcommand{\ill}{\textcolor{red!60!black}{[\,\ldots\,]}}

% Lecture douteuse : le mot est donné, et signalé comme incertain.
\newcommand{\uncertain}[1]{\uline{#1}}

% Ajout de l'éditeur — une lettre manquante, un mot sous-entendu.
\newcommand{\add}[1]{\textcolor{blue!50!black}{[#1]}}

% Biffé par Grothendieck lui-même. Ce qu'il a rayé fait partie du document :
% une rature dit souvent où la pensée a changé de direction.
\newcommand{\struck}[1]{\sout{#1}}

% Note du transcripteur, sur l'état matériel du feuillet.
\newcommand{\note}[1]{\par\smallskip{\small\color{gray!60!black}\textit{[#1]}}\par\smallskip}

% Note marginale de Grothendieck, distincte du corps du texte.
\newcommand{\marginal}[1]{\par\smallskip\noindent\hspace{1em}%
  {\small\color{gray!60!black}#1}\par\smallskip}

% --- Titre ------------------------------------------------------------------

\AtBeginDocument{%
  \begin{center}
    {\large\bfseries Fonds Alexandre Grothendieck --- cote n\textsuperscript{o}~\grothFolder}\\[2pt]
    {\itshape\grothFoldertitle}\\[4pt]
    {\small feuillets \grothFirst--\grothLast\ (lot \grothBatch)}\\[2pt]
    {\footnotesize\color{gray!60!black}Transcription établie d'après le fac-similé numérisé
     par l'Université de Montpellier.\\
     Les lectures incertaines sont soulignées, les illisibles notées [\,\ldots\,],
     les ajouts entre crochets.}
  \end{center}
  \vspace{1em}}

\endinput


\folder{154}
\batch{6}
\pages{101}{120}
\dating{1983-1984}
\watermark{Édition de démonstration}
\foldertitle{[Système de pseudo-droites] : notes manuscrites (1983-1984, s.d.)}

\begin{document}

\page{101}
\note{page de dessins en couleur. À gauche, une carte en forme de disque : des faces pentagonales et hexagonales bordées de brun, des arcs orange et vert clair au bord, et le graphe dual tracé en bleu (sommets bleus au centre des faces, arêtes bleues qui traversent les arêtes brunes) ; des tirets en marge prolongent certaines arêtes. En haut à droite, au crayon, un cercle traversé de droites. Au milieu et en bas, deux droites horizontales vert clair coupées de droites orange, avec une courbe vert clair qui serpente autour d'elles (un balayage). Entre les deux, une rangée de petits triangles au crayon, certains hachurés, au-dessus et au-dessous d'une droite, avec de petites légendes au crayon, en partie lisibles : « \uncertain{pas de tr.} », « \uncertain{s-g}, $g$ », « \uncertain{pas de tr.} $s$-$g$ », et plus bas « \uncertain{à g.} ou à \uncertain{droite} » ; deux légendes à gauche sont biffées}

\page{102}
\note{page de dessins en couleur : plusieurs disques bordés de vert clair traversés de pseudo-droites orange, certains au crayon, quelques secteurs hachurés ou coloriés ; un petit disque partagé par un arc brun ; un demi-disque et un segment de disque traversés d'arcs orange. En bas à droite, un grand disque traversé de pseudo-droites orange et d'un faisceau de courbes vert clair qui le parcourent de haut en bas ; les points d'intersection sont marqués, et le bord porte, au crayon, des numéros (lus $1$, $2$, $8$, $9$, $11$, $15$, $16$, $18$, en partie incertains)}

\page{103}
\note{en marge gauche, trois petits dessins en couleur : deux paires de droites orange qui se croisent sur une droite vert clair, celle-ci contournant l'un des croisements (vert foncé) ; les mêmes, avec des flèches d'orientation $\rightarrow$, $\leftarrow$ et de petites flèches verticales ; un croisement seul avec deux flèches}
\uncertain{Opération} de \uncertain{rotation} autour d'une \uncertain{face}

par une flèche de transition

$(D, \omega, a)$, $D$ position relative, $a$ arc de transition \add{orientée} sur $\widetilde{D}$, on trouve $(D', \omega', a')$ où \struck{on balaye sur l'arc de transition,}
\note{« orientée » est écrit au-dessus de la ligne, entre « relative » et « $a$ » ; sous le $D$ du triplet, un petit trait souligné, sans doute une marque ; « où » est ajouté au-dessus du passage biffé}
\begin{itemize}
\item[a)] $D'$ déduit de $D$ en \uncertain{transitant} à \uncertain{travers} $a$
\item[b)] $\omega'$ est l'orientation \emph{opposée} à celle déduite de $\omega$ par transport
\item[c)] $a'$ est l'arc de transition sur $\widetilde{D}'$, qui \emph{suit} l'arc de transition transporté de $a$, dans la direction $\omega'$ ;
\end{itemize}

\emph{NB} Le balayage de l'arc de spécialisation [croquis] donne [croquis].

Le balayage de l'arc de spécialisation [croquis] donne [croquis].

L'antibalayage de [croquis] donne [croquis].
\note{chaque « [croquis] » est un petit dessin en couleur dans la ligne : un croisement de deux droites orange sur la droite vert clair, l'arc de spécialisation noirci et fléché ; après « donne », la droite verte contourne le croisement par-dessus (bosse), l'arc noirci passé de l'autre côté. Dans la troisième ligne, les flèches sont dirigées vers la gauche}

Pour $\rho(D, \omega, a) = (D', \omega', a')$, \struck{\ill{}} \add{huit} cas à considérer, se groupant par deux paires, correspondant respectivement à $a$ est un arc de spécialisation, un arc de générisation.
\note{« huit » est écrit au-dessus du mot biffé ; notre lecture de ce mot ajouté est incertaine}

\note{au-dessous, un tableau de dessins en couleur, en deux groupes à accolade marqués $\mathrm{I}$ et $\mathrm{II}$, chacun de deux lignes numérotées 1) et 2) : à gauche, deux paires de droites orange qui se croisent sur la droite vert clair $D$, orientée par $\omega$ (flèche vers la gauche), avec l'arc $a$ marqué au premier croisement et, en pointillé, l'arc $a'_0$ au second ; puis « donne », et à droite la position $D'$, orientée par $\omega'$ (flèche vers la droite en $\mathrm{I}$, vers la droite ou la gauche en $\mathrm{II}$), qui passe au-dessus du premier croisement (groupe $\mathrm{I}$) ou au-dessous (groupe $\mathrm{II}$), l'arc $b$ marqué en pointillé au premier croisement et l'arc $a'$ au second. Dans $\mathrm{I}$ 2), la seconde paire de droites est surchargée d'un trait sombre}

Dans tous les cas, pour déterminer $a'$, on peut soit prendre l'arc $a'_0$ \add{de transition} qui suit $a$ sur $\widetilde{D}$ dans la direction \emph{opposée} à $\omega$, et prendre son \struck{\ill{}} \add{transporté} par balayage \add{ou antibalayage} ; soit transporter d'abord $a$ en $b$ par la dite \uncertain{spécialisation}, et prendre le successeur de $b$ pour $\omega'$, parmi les arcs de transition sur $\widetilde{D}'$. Cela résulte du fait que dans une opération de spécialisation, l'isom.\ can.\ entre $\operatorname{Pol}(D)$ et $\operatorname{Pol}(D')$ induit une bijection entre l'ens.\ des arcs de transition sur le premier,
\note{colonne de droite de la page ; « de transition », « transporté » et « ou antibalayage » sont écrits au-dessus de la ligne}

\page{104}
\emph{non adjacents} à \emph{$a$}, sur l'un des arcs de transition \struck{\ill{}} \struck{sur} \add{de} $\operatorname{Pol}(D')$ distincts de l'antipodique de $b$.
\note{la phrase continue celle du bas de la page 103}

\note{suit une rangée de cinq dessins en couleur reliés par des flèches $\mapsto$ : un cycle de positions relatives d'une droite vert clair par rapport à deux paires de droites orange et à un faisceau de droites parallèles marquées $\lambda$, $\alpha$, $\beta$, $\gamma$, $\mu$ ; la droite verte contourne tour à tour l'un ou l'autre croisement, l'arc $a$ et, en pointillé, l'arc $b$ sont marqués, et des flèches $\leftarrow$, $\rightarrow$ donnent l'orientation. Les dessins sont étiquetés $(1)$, $(2)$, $(3)$, $(4)$, $(1)$, le dernier revenant au premier ; les secteurs hachurés au crayon portent des numéros entourés $1$, $2$, $1'$, $2'$, $3'$, $4'$, et les mentions « $3' = 1$ » et « $5' = 1$ » ; sous les deux premiers, des chiffres romains surchargés}

\emph{NB} Pour chaque position $(D, \omega, a)$ dans le cycle, on a indiqué en pointillé l'arc de transition $b$ qui suit $a$ suivant l'orientation $\omega$.

Il faudrait regarder à part les cas où on part avec un arc de \struck{transition} \add{spécialisation} exceptionnel, correspondant à un balayage de couloir. Il ne devrait pas y avoir de différence.

\marginal{arcs de spécialisation, arcs de générisation $\Big\rbrace$ arcs de transition ; \emph{arcs \struck{de \ill{}} \add{d'isotopie}}}
\note{à droite, les deux premières lignes sont réunies par une accolade vers « arcs de transition » ; la troisième est soulignée, « d'isotopie » écrit au-dessus du mot biffé}

\[
2\nu - 2 + (2 + 2) = 2\nu + 2
\]
dont
\note{le premier signe peut aussi se lire « $\geqslant$ » ; la ligne s'arrête sur « dont »}

\note{la moitié inférieure de la page est faite de dessins en couleur : deux paires de droites orange coupées par une droite vert clair, les sommets $x$, $y$, $z'$, $y'$ reliés par des chemins en tirets vert foncé ; le même avec des points $z$, $t$ et un chemin en tirets ; deux rectangles bordés de vert clair d'où partent des rayons en étoile, les côtés marqués $\alpha$, $\beta$, $\gamma$ et les coins $x$, $y$, $x'$, $y'$, étiquetés $(1)$, $(2)$, $(3)$, $(4)$, entourés de petits carrés vert clair ; un petit carré vert clair avec des flèches $\lambda$, $\mu$ entre $x$, $y$, $x'$ ; un faisceau de droites orange presque parallèles, la droite verte passant par-dessus ; enfin, en bas à droite, une droite vert foncé coupée de $\nu$ droites orange (accolade $\nu$, sous-groupe marqué $\nu - 2$), les deux extrémités marquées $2$}

\page{105}
\note{page presque vide : au crayon, à gauche, quatre cercles traversés de courbes, disposés en trèfle, et un cinquième au-dessous ; ailleurs, des traces très pâles de croquis (droites qui se croisent, petits arcs en tirets)}

\page{106}
\note{dessins en couleur : deux paires de droites orange qui se croisent, entre elles un faisceau de droites fléchées $\lambda$, $\alpha$, $\beta$, $\gamma$, $\mu$ ; une droite vert clair et une droite en tirets, l'arc $x_0 y$ noirci ; à côté, la même configuration après une spécialisation, la droite verte passant au-dessus du premier croisement, avec les marques $x'_1 = x_0$, $y'_1 = x'$, $y'$, $z'$ et les flèches $-\lambda$, $-\alpha$, $-\beta$, $-\gamma$, $-\mu$. Plus bas, un croisement avec une bande hachurée qui s'enroule, les points $x_0$, $x$, $y$, $y'$, $z'$ ; un petit repère avec les flèches $\lambda$, $\alpha$, $\beta$, $\gamma$, $\mu$ entre $x$, $y$ et $y'$ ; et plusieurs croisements traversés de droites vert clair et de tirets bleus, avec des hachures au crayon}

région générisée $\mathrm{I}$

$\to$ arc de spécialisation pour $\mathrm{I}$, de générisation pour $\mathrm{II}$

\emph{d'attache de $\mathrm{I}$ et $\mathrm{II}$}

région spécialisée $\mathrm{II}$
\note{légendes du deuxième dessin ; « d'attache de I et II » est surligné}

pointillé bleu (s'arrêtant à l'arc de \uncertain{second} spécialisation-générisation) est celui du découpage

gén $\to$ spéc.

spéc $\to$ gén ou gén-spéc\uncertain{ial.}\ suivant qu'on rencontre un triangle, ou un sommet
\note{« suivant qu'on rencontre un » est écrit sur deux lignes, et une accolade réunit « triangle » et « ou un sommet »}

\page{107}
\note{page presque vide : en haut à gauche, deux petits dessins en couleur, reprise des deux premiers de la page 106 : deux paires de droites orange, un faisceau de droites fléchées $\lambda$, $\alpha$, $\beta$, $\gamma$, l'arc $x_0 x_1$ noirci, et, à droite, la droite vert clair passant au-dessus du premier croisement, l'arc $x_0 x_1$ marqué au-dessous}

\page{108}
\note{page presque vide : au crayon, deux paires de droites qui se croisent et une courbe qui les contourne ; ailleurs, des traces pâles de petits arcs en tirets}

\page{109}
\note{dessins en couleur sans texte : des droites vert clair orientées (flèche vers la gauche) coupées de droites rouges, verticales ou croisées deux à deux ; les droites vertes montent en marche d'escalier ou contournent les croisements, et des arcs sont renforcés en noir ; de fines flèches au crayon, numérotées $1$, $2$, $3$, indiquent les déplacements successifs. En bas à gauche, un croisement de deux droites rouges sur deux droites (vert clair et vert foncé), avec un petit cercle fléché autour du point}

\page{110}
\note{feuille en hauteur. En haut à gauche, un grand dessin en couleur de la bande supersingulière : deux droites horizontales vert clair, une ligne médiane vert foncé et en tirets, des droites rouges qui s'y croisent en sabliers et en losanges, et deux grands arcs rouges en haut et en bas ; légendes « $n - 1$ », « $\nu' - 1$ », « $\nu - 1$ », « $(n-1) - (\nu - 1) - (\nu' - 1)$ », et, à droite, « $n - \nu$ », « $(\nu - 1)$ », « $\nu$ »}

\struck{$n - \nu - \nu' + 1$}
\[
n - (\nu + \nu' - 1) = (n-1) - (\nu - 1) - (\nu' - 1)
\]
\[
n - \nu - (\nu' - 1)
\]
\note{à gauche, « $n - 1$ » deux fois, la seconde biffée ; au-dessous, un rectangle vert clair dont les grands côtés, renforcés en brun et coupés de petits traits rouges, sont marqués $n - \nu$, et les petits côtés $(\nu - 1)$, à droite « $(\nu - 1) + 1$ »}

\note{au milieu de la page, dessins en couleur : un rectangle vert clair dont les côtés portent de petits traits rouges, marqués $n - \nu$ (grands côtés) et $\nu - 1$ (petits côtés), avec un « $2(\nu - 1)$ » entouré à côté ; un faisceau de droites rouges en étoile ; un rectangle $a\,b\,a'\,b'$ traversé de droites rouges qui se croisent, avec une ligne médiane en tirets ; à droite, une droite vert clair qui passe d'une hauteur à l'autre à travers un croisement de droites rouges, avec les points $a$, $a'$, $b$, $b'$ ; un rectangle vert clair aux coins marqués de petits carrés}

\[
n - \nu \qquad n - \nu - (\nu' - 1) \qquad n - \nu \qquad n - \nu' \qquad n - \nu' \qquad n - 1 \qquad \nu
\]
\[
n - \nu \qquad 2n - \nu - \nu' + 2 \qquad 2n - (\nu + \nu' - 2) \quad \underline{\varepsilon \geqslant 2}
\]
\[
2n - 2(\nu + \nu' - 1) \qquad 2n - (\nu + \nu' - 1)
\]
\struck{$2\nu$} $2n - (\nu + \nu' - 1)$
\note{calculs épars, notés dans l'ordre de la page, de haut en bas ; dans « $2n - \nu - \nu' + 2$ », les deux $\nu$ sont barrés d'un trait oblique ; le $n$ de « $2n - (\nu + \nu' - 2)$ » est surchargé}

\note{en bas à gauche, quatre dessins en couleur de demi-cercles ou d'arcs vert clair aux extrémités marquées, avec des accolades au crayon portant des signes : « $-1$ » en haut et en bas, « $+1$ » sur les côtés et à l'intérieur dans le premier ; « $-1$ », « $+1$ » dans le deuxième ; dans les deux derniers, les côtés marqués « $\varepsilon$ » et l'intérieur « $+1$ », « $-1$ »}

\[
\varepsilon = \begin{cases} +1 & \text{si la face spéciale est du type sommet (ultra-\ldots)} \\ -1 & \text{sinon} \end{cases}
\]
\note{après « ultra- », un mot biffé, peut-être « spécial » ; « face » et « spéciale » sont d'une lecture incertaine}

\uncertain{Comment} \ill{} trouver la \uncertain{valeur} de $\varepsilon$ ?

\page{111}
\note{dessins en couleur : deux droites vert clair $D$ et $\overline{D}$ qui se croisent en un point multiple d'un faisceau de droites orange, de part et d'autre d'une droite brune horizontale, avec deux flèches courbes qui indiquent le passage de l'une à l'autre ; sur les droites orange, de chaque côté, des points marqués $u$, $v$, $w$, $u'$, $v'$, $w'$ et leurs barrés $\bar{u}$, $\bar{v}$, $\bar{w}$, $\bar{u}'$, \ldots, au milieu $r$, $s$, $t$, $\bar{r}$, $\bar{s}$, $\bar{t}$, à droite $a$, $a'$, $\bar{a}$, $\bar{a}'$ et $x$, $y$, $z$, $x'$, $y'$, $z'$ avec leurs barrés ; en bas à gauche une étiquette « $\overline{D}{}^{\ast}$ » (\uncertain{lue}). Au-dessus, une première version du même dessin, moins complète. Plus bas, un croisement de droites orange enjambé par une courbe vert clair en tirets, avec des flèches et deux arcs (noir et bleu) marqués ; à droite, deux petits faisceaux orange sur une droite vert clair}

\emph{NB} $x$ et $\bar{x}$ (ou encore $x'$ et $\bar{x}'$) \uncertain{ne} \uncertain{sont} pas sur la \uncertain{même} droite !

\page{112}
\note{à gauche, deux dessins en couleur : deux droites orange $D_1$, $D_2$ qui se croisent en $s$, une courbe vert foncé qui les enjambe au-dessus d'une droite vert clair et d'une droite en tirets, des hachures et des flèches au crayon ; étiquettes $D_1$, $D_2$, en bas « $D_2$ », « $D_4$ » (\uncertain{lus}), et dans le second $D$, $D'$, $D'_1$, et une accolade marquée $\Delta_D$ ; en marge, « $\Delta_{D'}$ »}

\emph{non} ! (le sommet $s$ disparaîtrait !) Et $D_1 \cap \Delta_{D'}$ \struck{\ill{}} \uncertain{va} sur $D_2 \cap \Delta_D$ et vice-versa. [Pourtant, ça vaut la peine de regarder ce que ça donnerait !)
\note{au-dessus du mot biffé, un mot illisible ; le crochet ouvrant se ferme sur une parenthèse}

Ici, le sommet est conservé, \struck{et dans} \struck{les droites} $D_1 \cap \Delta_D$ se continue en $D_1 \cap \Delta_{D'}$ \uncertain{c'est égal} --- mais dans le carré \uncertain{obtenu} \uncertain{par} un recollement \uncertain{là}, je n'arrive pas à suivre les « circuits autour d'un trou »

\note{à droite, deux dessins en couleur : un croisement de droites orange enjambé par une courbe vert foncé, au-dessus d'une droite vert clair et de tirets, avec une flèche verticale ; puis la même configuration où la droite vert foncé passe par le croisement}
\uncertain{Ici} \uncertain{aussi} $s$ disparaît au voisinage de l'arc de recollement

\note{en bas, deux disques : l'un bordé de vert foncé, traversé de deux pseudo-droites orange qui se croisent près du bord en haut et en bas, une petite croix orange dans un secteur ; l'autre, bordé de vert clair, déformé en pomme, les deux pseudo-droites confondues au milieu}

\page{113}
\note{à gauche, au crayon, deux figures en huit (deux cercles tangents), traversés de courbes et de petits traits, et, en bas, un cercle avec des points marqués $u$, $v$, $x^{\ast}$, $x$, $x + 1$ et un « $-1$ » entouré. Au milieu, des arcs en couleur (vert clair, traits rouges et bleus), les extrémités marquées $u$, $v$, $u'$, $v'$, les croisements $\varepsilon'$, $\varepsilon''$, les accolades $\alpha'$, $\alpha''$, $\gamma$, et « $\alpha' + \alpha'' + \gamma$ » ; des flèches doubles $\Updownarrow$ les relient deux à deux. Un « $N$ » isolé à gauche}

$\alpha', \alpha'' \geqslant 1$, $\gamma \geqslant 2$ \qquad $\alpha' + \alpha'' + \gamma \geqslant 4$
\note{le « $1$ » est surchargé ; dans la seconde formule, le $\gamma$ est surchargé}

$N$ nb de « sommets verts » \add{de $F$} (i.e.\ sommets \struck{rouges} \add{d'ordre $4$})

$N'$ --------- de $F'$

\[
N' = N - 1
\]
\note{au-dessous, au crayon, un cercle traversé de courbes}
\[
N' = N
\]
\note{au-dessous, au crayon, un huit (deux cercles tangents) traversé de courbes, et à côté le même huit en couleur, vert clair et vert foncé, des arcs renforcés en brun et marqués de petits traits rouges. Un « $N$ » isolé en haut à droite de la page}

\page{114}
\note{page de dessins : en haut à gauche, un hexagone et un pentagone vert clair accolés, des croix rouges sur les arêtes ; au crayon, deux cercles tangents traversés de courbes, le point de contact marqué, avec les accolades $\alpha_0$, $\alpha_1$, $\alpha_2$, $\alpha''$ et « $n - \nu$ » ; en haut à droite, un croquis au crayon d'un pentagone et d'une droite coupée de droites, avec un arc vert clair au-dessous. En bas, un huit au crayon marqué « $N' = N - 1$ », des arcs vert clair aux points marqués « $1$ », « $2$ », « $3$ » (avec des signes $\sim$), des cercles au crayon, et à droite un huit en couleur (vert clair, vert foncé, arcs renforcés en brun) aux arcs marqués « $1$ », « $1{+}$ », « $3$ », « $13$ »}

\page{115}
\note{page de dessins en couleur : sept demi-cercles vert clair aux extrémités marquées, certains arcs renforcés en brun ; dans trois d'entre eux, des droites rouges qui descendent du haut, verticales ou obliques, et de petits traits rouges sur le bord ; dans les trois de droite, le sommet marqué « $+1$ », les côtés marqués $\varepsilon_1$, $\varepsilon_2$, et des segments intérieurs au crayon portant $\varepsilon_1$, $\varepsilon_2$ et « $-1$ » ; en haut à droite, un cercle au crayon}

\page{116}
\note{sur une feuille de listing imprimée (bannière de travail), des dessins en couleur : à droite, une droite vert clair coupée de quatre paires de droites orange, les arcs $a$, $b$, $b'$ marqués, des hachures ; en bas à gauche, un carré vert clair (côtés marqués $(1)$, $(2)$, $(3)$) aux coins duquel partent des arêtes vert foncé et des arcs orange marqués $\lambda$, $\mu$, $\lambda'$, $\mu'$ ; au milieu, un croquis au crayon d'un triangle avec les mêmes lettres ; à droite, le même carré, ses quatre côtés marqués $\mathrm{I}$, $\mathrm{II}$, $\mathrm{III}$, $\mathrm{IV}$, rempli d'arcs orange et de flèches, avec les lettres $a$, $b$, $a'$, $b'$, $\lambda$, $\mu$, $x$ et trois points d'interrogation}

\page{117}
\note{toute la page est barrée d'un grand X tracé en deux traits verts}
\uncertain{donc} on trouve (\ill{} qui est dans le \ill{} \ill{}) les \uncertain{trois} \ill{}
\note{les deux premières lignes continuent une phrase commencée hors de cette page ; seuls quelques mots se lisent}

\note{deux dessins en couleur : un « couloir » rectangulaire (bords vert clair et vert foncé, tirets verts) traversé de droites rouges, un arc rouge le fermant en haut, avec deux flèches vers la légende « au choix » ; même dessin, variante. À côté du premier, « \struck{\ill{} \ill{}} ou »}

\uncertain{d)} \struck{On trouve} on \uncertain{voit} encore, $\mathcal{E}$, un \ill{} \struck{\ill{}} couloir \ill{} \uncertain{canonique} autour de $K_0$ comme ceci (qu'il \uncertain{avait} utilisé pour la construction de $\mathcal{G}$ au voisinage de $K_0$) et on remplace $F$ par $\mathring{F}$ \uncertain{déduit} en \uncertain{enlevant} $F \cap \mathring{C}$.
\note{la lettre de l'alinéa, écrite sous le biffage, se lit « c) » ou « d) » ; les signes lus $\mathcal{E}$ et $\mathcal{G}$ sont incertains}

\note{dessin en couleur : une bande horizontale (droites vert clair, ligne médiane brune et verte, tirets jaunes) traversée de droites rouges verticales et de losanges rouges, fermée en haut et en bas par des arcs rouges}

\emph{NB} Ce couloir est isomorphe (dans la catégorie topologique convenable) au voisinage tubulaire de $\widetilde{D}_i$ dans $\widetilde{X}$

e) Pour les faces exceptionnelles \add{$F^{\ast}$, \uncertain{adjacentes} à $D_i$}, on retrouve le système $\Sigma$, \struck{\ill{}} \add{\uncertain{affaibli}} par le choix de $D_i$, en prenant « la partie de $F^{\ast}$ au-dessus de » la bande principale $\delta$, et faisant les modifications d'usage au voisinage des pts critiques.

\note{dessin en couleur : un hexagone bordé de vert clair et de noir, rempli de courbes rouges qui s'entrecroisent}

f) Attention, à cause de la rectification b), il n'est plus correct de caractériser les droites exceptionnelles \add{de $K$} comme étant celles qui ne contiennent pas de sommets de $L'$ (si elles sont exceptionnelles au sens d'une telle propriété, c'est plutôt qu'elles \emph{peuvent} en contenir plusieurs !). \struck{Rectification on peut} \add{Mais ici \struck{tout}} \uncertain{construisons} $K'_0$ comme la réunion \struck{plus} \struck{constituées d'arcs des faces exceptionnelles} \struck{isolées plus} \struck{un \ill{}} \uncertain{des} composantes connexes du \struck{$K$} \uncertain{déploiement} $K'$ qui n'ont pas de bord \emph{global}, chose maintenant \uncertain{tout} naturel pour introduire ensuite la donnée d\ldots
\note{la page s'arrête sur « d’- », coupé en fin de ligne ; les trois lignes biffées sont d'une lecture incertaine}

\page{118}
\note{en haut à gauche, un dessin en couleur de la bande supersingulière (droites vert clair, ligne médiane en tirets rouges et bleus, droites rouges en losanges, arcs rouges en haut et en bas) ; en haut à droite, deux arcs vert clair en forme d'arche, aux sommets marqués, sous la légende « \struck{\ill{}} $\alpha \geqslant 2$ sommets verts clairs » ; ailleurs, des croquis au crayon : des droites $D$, $D'$, $D''$ qui se coupent, les points $a$, $a'$, $x$, $x'$ et les étiquettes $\omega_D$, $\omega_{D'}$, $\omega_{D''}$, un repère fléché}

\[
(\mathcal{X}, K, (\delta_s), (\nu_s), (\nu_{\ill{}}))
\]
\note{sous le troisième terme, un mot illisible ; notre lecture des lettres $\delta$ et $\nu$ est incertaine}

\[
\begin{array}{ccc}
\omega(D) & \xrightarrow[\text{comb.\ } \beta]{\text{rot } \alpha} & \omega(D') \\
\| & & \| \\
\omega(\mathcal{X}, D) & & \omega(\mathcal{X}, D')
\end{array}
\]
\note{sous « comb. $\beta$ » : « via les $\operatorname{Pol}(D')$ » ; « rot » est \uncertain{lu} ; en face, « $\alpha = -1$ »}

$N$ nb de \uncertain{faces}

\uncertain{1)} $N$ pair, $\alpha = \beta$, \ill{}

\[
\omega_D \wedge \omega_{D'} \wedge \lbrace D, D' \rbrace \qquad \omega_{D'} \wedge \omega_{D''} \wedge \lbrace D', D'' \rbrace
\]
\note{dans la première formule, le premier terme de l'accolade est surchargé}
\[
\rho_D \simeq \omega_{D'} \qquad \rho_{D'} \simeq \omega_D \qquad \rho_D \wedge \omega_{D'} \qquad \rho_{D'} \wedge \omega_D
\]
\note{ailleurs sur la page, isolés : « $\omega_D \wedge \omega_{D'}$ », « $\omega_{D'}$ », « $\omega_D$ », « $\rho$ », et, près d'un croquis de $D$ et $D'$ qui se coupent, « $(a, a') = x$ »}

\page{119}
\note{à gauche, des dessins en couleur et au crayon : une droite vert clair traversée par deux croisements de droites rouges, entourée d'un contour vert foncé fléché dont les arcs sont numérotés $1$ à $4$ ; une courbe qui passe d'une hauteur à l'autre à travers un croisement, les régions hachurées ; un croisement enjambé par des courbes, régions numérotées $1$ à $7$ ; un disque et un ruban en huit au crayon ; en bas au milieu, un croisement de droites rouges avec deux courbes vert foncé et vert clair}

\[
(n - 2) + (n - 3) = 2n - 5 \qquad 2(n-1)(2n-5) \qquad 2n
\]

\[
\chi(\mathcal{X}) = s_0 - s_1 + s_2 = 1 - g
\]
Les sommets sont d'ordre $4$

On trouve donc
\[
2 s_1 = \underbrace{4 s_0}_{\sum_{s \in S} \nu_s} = 2 \sum_{f \in F} n_f
\]
\struck{\ill{}} donc $s_1 = 2 s_0$

donc $\chi(\mathcal{X}) = s_2 - s_0$ \qquad $g = 1 + s_0 - s_2$

où $s_2 = {}$nombre de p.r.\ (les supersingulières exclues)

$s_0 = {}$nombre de sommets i.e.\ \struck{des} des couples $(D, a)$, où $D$ est une p.r.\ et $a$ un arc \add{d'\uncertain{isotopie}} \struck{\ill{}} \struck{\ill{} \ill{} \ill{}} sur $D$ du type gén-gén.

b) des \struck{\ill{} \ill{}} \add{\uncertain{demi}-\ill{} du type} \struck{de $\Sigma$} $(0, 1)$ de $\Sigma$ ($= 2n(n-1)$ dans \uncertain{les} \uncertain{cas} \uncertain{stables})
\note{« d'isotopie » est ajouté au-dessus de la ligne ; « b) » n'est précédé d'aucun « a) » sur la page}

\page{120}
\note{sur une feuille de listing imprimée (bannière de travail), des croquis : en haut, en couleur, un ruban vert foncé fermé autour d'une droite noircie, traversé par deux paires de droites rouges ; ailleurs, au crayon, des droites qui se croisent deux à deux sur une droite fléchée, un parallélépipède en perspective, des quadrilatères aux sommets marqués d'un point et des ovales traversés de droites}

\end{document}
